(1) 无涡旋区满足$\boldsymbol B-\lambda_L^2\nabla^2\boldsymbol B=0$；涡旋相位绕数导致$(1-\lambda_L^2\nabla_\perp^2)B_z=\Phi_0\delta^{(2)}(\boldsymbol r)$。二维修正Helmholtz算符的Green函数为$K_0(r/\lambda_L)/(2\pi\lambda_L^2)$，因此对$r\gtrsim\xi$，$B_z(r)=\Phi_0 K_0(r/\lambda_L)/(2\pi\lambda_L^2)$，远离涡旋按$e^{-r/\lambda_L}/\sqrt r$衰减。
(2) 极坐标积分$\Phi=\int_0^\infty B_z\,2\pi r\,dr=\Phi_0\int_0^\infty uK_0(u)\,du=\Phi_0$，其中$\int_0^\infty uK_0(u)\,du=1$，也可把源方程全平面积分得到同样结论。Ampere定律给$J_\phi=-\partial_rB_z/\mu_0=\Phi_0 K_1(r/\lambda_L)/(2\pi\mu_0\lambda_L^3)$，电流方向与涡旋通量取向一致。
(3) London单位长度自由能$F_\ell=(2\mu_0)^{-1}\int_\xi^\infty[B_z^2+\lambda_L^2(\partial_rB_z)^2]\,2\pi r\,dr+E_{\rm core}$。在无源区域用$B-\lambda_L^2\nabla^2B=0$并分部积分，面积项化为核心圆边界；对$\xi\ll r\ll\lambda_L$，$K_0(r/\lambda_L)\simeq\ln(\lambda_L/r)+O(1)$，遂有$F_\ell=\Phi_0^2\ln(\lambda_L/\xi)/(4\pi\mu_0\lambda_L^2)+O(\Phi_0^2/(\mu_0\lambda_L^2))$。对数系数无需知道芯内微观细节。
(4) 单位长度涡旋在外磁场$H$中的Gibbs能量差为$\Delta G_\ell\simeq F_\ell-\Phi_0H$，因为磁场源对新增通量作功。令$\Delta G_\ell=0$，得到$H_{c1}\simeq\Phi_0\ln\kappa/(4\pi\mu_0\lambda_L^2)$，对应$B_{c1}=\mu_0H_{c1}$。核芯能改变量级为对数项以外的常数；接近$H_{c1}$的涡旋间相互作用、退磁场和边界势垒都不能仅由孤立涡旋模型决定。
